% !TeX root = ../../ONMIPA Alin.tex
% Acuan: Zhang (2009), Bab 5, Soal 5.33, 5.36--5.51;
% Zhang (2011), bagian 1.4, hubungan kernel--citra dan adjoin.
% Seluruh operator dibahas dalam dimensi hingga; tidak memakai SVD, QR, atau LU.
\chapter[Operator pada Ruang Hasil Kali Dalam]{Operator pada Ruang\\Hasil Kali Dalam}
\label{ch:operator-hkd}

Hasil kali dalam memungkinkan sebuah transformasi dipindahkan
dari satu argumen ke argumen lainnya melalui operator adjoin.
Konsep ini menghubungkan kernel dengan komplemen ortogonal citra,
serta menjelaskan sifat geometris operator uniter dan adjoin diri.
Semua ruang pada bab ini berdimensi hingga atas $\R$ atau $\C$.
Konvensi hasil kali dalam tetap linear pada argumen kedua.

\section{Representasi Fungsional Linear}

Fungsional linear adalah transformasi linear $\ell:V\to\F$.
Untuk setiap $z\in V$, aturan $x\mapsto\langle z,x\rangle$
merupakan fungsional linear. Dalam dimensi hingga, semua
fungsional linear dapat diperoleh dengan cara tersebut.

\begin{teorema}[][Representasi Riesz dalam dimensi hingga][op:riesz]
    Untuk setiap fungsional linear $\ell:V\to\F$, terdapat
    tepat satu $z\in V$ dengan
    \[
        \ell(x)=\langle z,x\rangle\qquad\text{untuk semua }x\in V.
    \]
    Jika $V\ne\{0\}$, maka
    $\max_{\|x\|=1}|\ell(x)|=\|z\|$.
\end{teorema}
\begin{proof}
    Ambil basis ortonormal $e_1,\ldots,e_n$ dan letakkan
    $a_j=\ell(e_j)$. Untuk $x=\sum x_je_j$,
    $\ell(x)=\sum a_jx_j$.
    Vektor $z=\sum\overline{a_j}e_j$ memenuhi
    $\langle z,x\rangle=\sum a_jx_j$.
    Jika $z'$ juga merepresentasikan $\ell$, maka
    $\langle z-z',x\rangle=0$ untuk semua $x$.
    Pilih $x=z-z'$ untuk memperoleh $z=z'$.
    Cauchy--Schwarz memberi $|\ell(x)|\le\|z\|$ pada vektor satuan.
    Jika $z\ne0$, batas dicapai pada $x=z/\|z\|$; jika $z=0$,
    semua nilai fungsional nol.
\end{proof}

Pada ruang kompleks, pemetaan $z\mapsto\langle z,\cdot\rangle$
bersifat konjugat-linear terhadap $z$. Konjugasi pada koefisien
representer dalam bukti di atas tidak boleh dihapus.

\begin{exam}[Kendala linear sebagai hiperbidang ortogonal]
    Pada $\C^3$, ambil $\ell(x)=2x_1-ix_2+3x_3$.
    Tentukan representer, nilai maksimum $|\ell(x)|$ pada
    vektor satuan, dan solusi bernorma minimum dari $\ell(x)=b$.
\end{exam}
\begin{solution}
    Representernya $z=(2,i,3)^T$, dengan $\|z\|^2=14$.
    Maka maksimum yang diminta adalah $\sqrt{14}$ dan
    $\ker\ell=\operatorname{span}\{z\}^\perp$.
    Solusi minimum adalah $x_0=(b/14)z$ karena
    $\langle z,x_0\rangle=b$.
    Semua solusi lain berbentuk $x_0+w$ dengan $w\perp z$,
    sehingga norma kuadratnya $|b|^2/14+\|w\|^2$.
\end{solution}

\section{Adjoin dan Matriks Representasinya}

\begin{definisi}[][Adjoin transformasi linear][op:adjoin]
    Adjoin dari $T:V\to W$ adalah transformasi $T^*:W\to V$
    yang memenuhi
    \[
        \langle Tx,y\rangle_W=\langle x,T^*y\rangle_V
        \qquad(x\in V,\ y\in W).
    \]
\end{definisi}

\begin{teorema}[][Keberadaan dan keunikan adjoin][op:ada-adjoin]
    Setiap transformasi linear antara dua ruang hasil kali dalam
    berdimensi hingga mempunyai tepat satu adjoin.
    Jika matriks $T$ dalam basis ortonormal adalah $A$, matriks
    adjoin dalam arah sebaliknya adalah $A^*$.
\end{teorema}
\begin{proof}
    Pilih basis ortonormal pada kedua ruang, dan tulis koordinat
    $x,y$ sebagai $u,v$. Identitas
    \[
        (Au)^*v=u^*(A^*v)
    \]
    menunjukkan bahwa transformasi bermatriks $A^*$ memenuhi
    definisi adjoin. Untuk keunikan, jika $S$ juga memenuhi definisi,
    maka $\langle x,(S-T^*)y\rangle=0$ untuk semua $x,y$.
    Dengan memilih $x=(S-T^*)y$, diperoleh $(S-T^*)y=0$
    untuk semua $y$, sehingga $S=T^*$.
\end{proof}

Jika basis yang digunakan tidak ortonormal, misalkan matriks
hasil kali dalam domain dan kodomain adalah $G_V$ dan $G_W$.
Untuk $A=[T]_{\mathcal C\leftarrow\mathcal B}$, identitas adjoin memberi
\[
    [T^*]_{\mathcal B\leftarrow\mathcal C}=G_V^{-1}A^*G_W.
\]
Memang, kedua ruas identitas hasil kali dalam dalam koordinat
adalah $u^*A^*G_Wv$ dan $u^*G_V[T^*]v$.
Matriks adjoin baru sama dengan transpose konjugat matriks asal
ketika basis-basisnya ortonormal, atau dalam keadaan khusus yang sesuai.

\begin{exam}[Adjoin diri dalam hasil kali dalam berbobot]
    Pada $\R^2$ dengan $G=\diag(1,2)$, tinjau operator
    bermatriks $A=\begin{pmatrix}0&2\\1&0\end{pmatrix}$.
    Tentukan matriks adjoinnya.
\end{exam}
\begin{solution}
    Perhitungan memberi
    \[
        G^{-1}A^TG
        =\begin{pmatrix}1&0\\0&1/2\end{pmatrix}
         \begin{pmatrix}0&1\\2&0\end{pmatrix}
         \begin{pmatrix}1&0\\0&2\end{pmatrix}=A.
    \]
    Jadi operator sama dengan adjoinnya walaupun $A$ tidak
    simetris dalam basis standar tersebut.
    Vektor eigen $(\sqrt2,1)^T$ dan $(-\sqrt2,1)^T$
    ortogonal dalam metrik $G$, karena hasil kali dalamnya $-2+2=0$.
\end{solution}

\begin{prop}[Aljabar adjoin]
    Untuk transformasi dengan domain dan kodomain yang sesuai,
    \[
        (T^*)^*=T,\qquad
        (aS+bT)^*=\overline aS^*+\overline bT^*,\qquad
        (ST)^*=T^*S^*.
    \]
    Jika $T$ invertibel, maka $(T^{-1})^*=(T^*)^{-1}$.
\end{prop}
\begin{proof}
    Masing-masing rumus dapat diperiksa melalui identitas
    pendefinisi adjoin dan keunikannya.
    Misalnya,
    $\langle STx,y\rangle=\langle Tx,S^*y\rangle
    =\langle x,T^*S^*y\rangle$.
    Rumus invers mengikuti pengambilan adjoin pada
    $TT^{-1}=I$ dan $T^{-1}T=I$.
\end{proof}

\begin{exam}[Adjoin operator pada ruang matriks]
    Pada ruang matriks dengan hasil kali dalam Frobenius,
    tentukan adjoin $L(X)=AXB$, dengan ukuran matriks yang sesuai.
\end{exam}
\begin{solution}
    Sifat siklik trace memberi
    \[
        \langle AXB,Y\rangle_F
        =\tr(B^*X^*A^*Y)
        =\tr(X^*A^*YB^*)=\langle X,A^*YB^*\rangle_F.
    \]
    Jadi $L^*(Y)=A^*YB^*$.
    Secara khusus, adjoin dari $X\mapsto AX-XA$ adalah
    $Y\mapsto A^*Y-YA^*$.
\end{solution}

\section{Hubungan Kernel, Citra, dan Ortogonalitas}

\begin{teorema}[][Empat ruang fundamental melalui adjoin][op:empat-ruang]
    Untuk $T:V\to W$,
    \[
        \ker T^*=(\operatorname{im}T)^\perp,\qquad
        \operatorname{im}T^*=(\ker T)^\perp.
    \]
    Akibatnya,
    \[
        V=\ker T\oplus\operatorname{im}T^*,\qquad
        W=\operatorname{im}T\oplus\ker T^*,
    \]
    dengan kedua jumlah ortogonal, dan $\rank T=\rank T^*$.
\end{teorema}
\begin{proof}
    Vektor $y\in W$ berada di $(\operatorname{im}T)^\perp$
    tepat ketika $\langle Tx,y\rangle=0$ untuk semua $x$.
    Dengan identitas adjoin, hal ini ekuivalen dengan
    $\langle x,T^*y\rangle=0$ untuk semua $x$, yakni $T^*y=0$.
    Terapkan hasil ini pada $T^*$ dan ambil komplemen ortogonal
    untuk memperoleh identitas kedua.
    Dekomposisi mengikuti teorema komplemen ortogonal.
    Terakhir, $\rank T^*=\dim V-\dim\ker T=\rank T$.
\end{proof}

Selain itu,
\[
    \langle x,T^*Tx\rangle=\|Tx\|^2,
    \qquad \ker(T^*T)=\ker T.
\]
Satu inklusi kernel jelas; inklusi lainnya mengikuti bahwa
$T^*Tx=0$ mengakibatkan $\|Tx\|^2=0$.
Jadi $\rank(T^*T)=\rank T$.
Dalam masalah kuadrat terkecil, hubungan ini menjelaskan mengapa
matriks persamaan normal invertibel tepat ketika kolom matriks
asal bebas linear.

Jika $T$ surjektif, setiap persamaan $Tx=y$ mempunyai satu
solusi yang berada dalam $(\ker T)^\perp$; solusi tersebut
adalah solusi dengan norma minimum. Keunikannya mengikuti
dekomposisi ortogonal, sebab solusi lain berbeda sebesar
vektor dari $\ker T$.

\section{Isometri dan Operator Uniter}

\begin{definisi}[][Isometri linear dan operator uniter][op:isometri]
    Transformasi linear $T:V\to W$ disebut isometri jika
    $\langle Tx,Ty\rangle_W=\langle x,y\rangle_V$ untuk semua $x,y$.
    Isometri bijektif pada ruang yang sama disebut operator
    \emph{uniter} pada ruang kompleks dan \emph{ortogonal} pada ruang real.
\end{definisi}

\begin{teorema}[][Karakterisasi isometri linear][op:uji-isometri]
    Untuk transformasi linear $T:V\to W$, pernyataan berikut ekuivalen:
    \begin{enumerate}[label=(\roman*),nosep]
        \item $T$ mempertahankan hasil kali dalam;
        \item $\|Tx\|=\|x\|$ untuk setiap $x$;
        \item $T^*T=I_V$;
        \item citra suatu basis ortonormal $V$ adalah daftar ortonormal di $W$.
    \end{enumerate}
    Isometri selalu injektif. Jika $\dim V=\dim W$, ia bijektif
    dan $T^{-1}=T^*$.
\end{teorema}
\begin{proof}
    Pernyataan pertama memberi kedua dengan mengambil $y=x$.
    Arah balik memakai identitas polarisasi dan linearitas $T$.
    Identitas $\langle Tx,Ty\rangle=\langle x,T^*Ty\rangle$
    membuktikan ekuivalensi dengan pernyataan ketiga.
    Memeriksa pada pasangan anggota basis memberi pernyataan
    keempat; ekspansi koordinat membuktikan arah baliknya.
    Pelestarian norma memaksa kernel hanya nol.
    Jika dimensi sama, injektivitas memberi bijektivitas dan
    $T^*T=I$ kemudian memberi $T^*=T^{-1}$.
\end{proof}

\begin{exam}[Pelestarian norma tanpa linearitas]
    Pemetaan $C(z)=\overline z$ pada $\C$ mempertahankan norma
    dan jarak, tetapi bukan linear kompleks: $C(i)=-i\ne iC(1)$.
    Pemetaan translasi $x\mapsto x+a$ pada $\R^n$, dengan $a\ne0$,
    mempertahankan jarak tetapi tidak linear karena tidak
    memetakan nol ke nol. Jadi asumsi linearitas dalam teorema
    pelestarian norma tidak boleh dihilangkan.
\end{exam}

Dalam basis ortonormal, isometri mempunyai matriks $A$ dengan
$A^*A=I$. Jika matriks persegi, ini ekuivalen dengan uniter.
Pergantian antara dua basis ortonormal juga diberikan oleh
matriks uniter, sehingga representasi suatu operator berubah
melalui $A\mapsto U^*AU$.
Untuk basis dengan matriks metrik $G$, operator bermatriks $A$
mempertahankan hasil kali dalam tepat ketika $A^*GA=G$.

\section{Operator Adjoin Diri, Normal, dan Positif}

\begin{definisi}[][Kelas operator berdasarkan adjoin][op:kelas]
    Operator $T:V\to V$ disebut \emph{adjoin diri} jika $T^*=T$,
    \emph{antiadjoin diri} jika $T^*=-T$, dan \emph{normal}
    jika $T^*T=TT^*$.
    Operator adjoin diri disebut semidefinit positif jika
    $\langle x,Tx\rangle\ge0$ untuk semua $x$, dan positif definit
    jika ketaksamaan itu ketat untuk $x\ne0$.
\end{definisi}

Dalam basis ortonormal, sifat-sifat ini tepat sifat Hermitian,
anti-Hermitian, normal, serta positif pada matriks representasinya.
Operator adjoin diri dan operator uniter selalu normal.
Untuk operator adjoin diri, $\langle x,Tx\rangle$ real karena
sama dengan konjugatnya.

\begin{teorema}[][Teorema spektral untuk operator][op:spektral]
    Setiap operator adjoin diri mempunyai basis ortonormal
    yang terdiri atas vektor eigen, dengan nilai eigen real.
    Pada ruang kompleks, setiap operator normal mempunyai
    basis ortonormal vektor eigen.
\end{teorema}
\begin{proof}
    Pilih basis ortonormal awal. Matriks operator menjadi
    Hermitian pada kasus adjoin diri, atau normal pada kasus kedua.
    Terapkan teorema spektral matriks pada
    Teorema~\ref{eigen:teorema-spektral}.
    Kolom-kolom matriks perubahan basis uniter memberikan
    koordinat basis ortonormal vektor eigen operator.
    Pada ruang real dan operator adjoin diri, matriksnya simetris
    real sehingga basis eigen dapat dipilih real.
\end{proof}

Kerealan nilai eigen adjoin diri juga dapat dilihat langsung.
Jika $Tx=\lambda x$ dan $x\ne0$, maka
$\lambda\|x\|^2=\langle x,Tx\rangle$ real.
Jika $Tu=\lambda u$ dan $Tv=\mu v$ dengan $\lambda\ne\mu$,
identitas adjoin memberi $(\lambda-\mu)\langle u,v\rangle=0$;
jadi ruang eigen berbeda saling ortogonal.

\begin{exam}[Normal real belum tentu mempunyai basis eigen real]
    Matriks $R=\begin{pmatrix}0&-1\\1&0\end{pmatrix}$ ortogonal,
    sehingga normal. Namun, polinom karakteristiknya $t^2+1$
    tidak mempunyai akar real. Jadi tidak ada basis eigen real.
    Setelah ruang diperluas ke $\C^2$, nilai eigennya $i,-i$
    dan teorema spektral kompleks berlaku.
\end{exam}

Jika $T$ normal, ekspansi dalam basis eigen ortonormal memberi
$\|Tx\|=\|T^*x\|$ dan $\ker T=\ker T^*$.
Pada ruang real, kesamaan norma langsung mengikuti
$\langle x,(T^*T-TT^*)x\rangle=0$.
Oleh karena itu, operator normal memenuhi
$V=\ker T\oplus\operatorname{im}T$ secara ortogonal.
Operator normal yang nilpoten harus nol: pada ruang kompleks,
semua nilai eigen nol dan diagonalisasi memberi $T=0$;
untuk matriks real, argumen yang sama dapat diterapkan setelah
dipandang sebagai matriks kompleks.

\section{Hasil Bagi Rayleigh dan Dimensi Subruang Positif}

\begin{teorema}[][Batas hasil bagi Rayleigh][op:rayleigh]
    Misalkan $T$ adjoin diri pada ruang tak nol, dengan nilai
    eigen terkecil $\lambda_{\min}$ dan terbesar $\lambda_{\max}$.
    Untuk $x\ne0$,
    \[
        \lambda_{\min}\le
        \frac{\langle x,Tx\rangle}{\|x\|^2}
        \le\lambda_{\max}.
    \]
    Kesamaan pada masing-masing batas berlaku tepat ketika
    $x$ berada dalam ruang eigen nilai batas tersebut.
\end{teorema}
\begin{proof}
    Tulis $x=\sum c_je_j$ dalam basis eigen ortonormal.
    Maka hasil bagi yang dimaksud adalah
    \[
        \frac{\sum_j\lambda_j|c_j|^2}{\sum_j|c_j|^2},
    \]
    yakni rata-rata berbobot nilai-nilai eigen.
    Batas atas tercapai tepat ketika setiap koefisien untuk
    nilai eigen yang lebih kecil nol. Argumen batas bawah serupa.
\end{proof}

\begin{exam}[Maksimum bentuk kuadratik pada lingkaran satuan]
    Tentukan rentang $2x^2+2xy+2y^2$ untuk real $x^2+y^2=1$.
\end{exam}
\begin{solution}
    Matriks bentuk kuadratiknya $A=\begin{pmatrix}2&1\\1&2\end{pmatrix}$
    mempunyai nilai eigen $1$ dan $3$, masing-masing pada arah
    $(1,-1)^T$ dan $(1,1)^T$.
    Jadi minimum satu dicapai ketika $x=-y=\pm1/\sqrt2$,
    dan maksimum tiga ketika $x=y=\pm1/\sqrt2$.
    Semua nilai antara keduanya dicapai dengan mengubah
    koefisien dalam basis eigen ortonormal.
\end{solution}

\begin{teorema}[][Banyak nilai eigen positif][op:indeks-positif]
    Jika operator adjoin diri $T$ mempunyai tepat $r$ nilai
    eigen positif, dihitung dengan multiplisitas, maka $r$
    adalah dimensi terbesar subruang $U$ yang memenuhi
    \[
        \langle x,Tx\rangle>0\qquad\text{untuk setiap }x\in U\setminus\{0\}.
    \]
\end{teorema}
\begin{proof}
    Jumlah ruang eigen untuk nilai positif mempunyai dimensi $r$
    dan memenuhi ketaksamaan tersebut.
    Ruang $N$ yang dibangun oleh vektor eigen untuk nilai
    nonpositif berdimensi $n-r$.
    Jika $\dim U>r$, rumus dimensi memberi
    $\dim(U\cap N)\ge\dim U+(n-r)-n>0$.
    Vektor tak nol pada irisan menghasilkan nilai kuadrat
    positif menurut definisi $U$ dan nonpositif menurut definisi
    $N$, kontradiksi.
\end{proof}

Dengan mengganti $T$ oleh $-T$, diperoleh pernyataan serupa
untuk banyaknya nilai eigen negatif. Hasil ini sering mengubah
informasi tentang suatu subruang menjadi batas jumlah nilai
eigen tanpa menghitung polinom karakteristik.

\section{Latihan Pilihan}

\begin{latihan}\label{op:lat-riesz}
    Pada $\mathcal P_{\le2}(\R)$ dengan hasil kali dalam integral
    pada $[-1,1]$, tentukan representer dari $\ell(p)=p(0)$.
    Carilah konstanta terkecil $C$ sehingga
    \[
        |p(0)|^2\le C\int_{-1}^1|p(t)|^2\,dt
    \]
    untuk semua polinom pada ruang tersebut.
\end{latihan}

\begin{latihan}\label{op:lat-adjoin-berbobot}
    Pada $\R^2$, gunakan metrik $G=\begin{pmatrix}2&1\\1&2\end{pmatrix}$.
    Jika matriks $T$ dalam basis standar adalah
    $A=\begin{pmatrix}1&1\\0&2\end{pmatrix}$, tentukan matriks
    $T^*$. Apakah $T$ adjoin diri?
\end{latihan}

\begin{latihan}\label{op:lat-komutator}
    Untuk matriks Hermitian $A\in\Mat_n(\C)$, definisikan
    $D_A(X)=AX-XA$ pada ruang matriks dengan hasil kali dalam
    Frobenius. Buktikan bahwa $D_A$ adjoin diri.
    Buktikan pula bahwa $D_A$ semidefinit positif tepat ketika
    $A$ merupakan matriks skalar.
\end{latihan}

\begin{latihan}\label{op:lat-titik-tetap}
    Jika $U$ uniter, buktikan dekomposisi ortogonal
    \[
        V=\ker(U-I)\oplus\operatorname{im}(U-I).
    \]
    Tafsirkan kedua subruang sebagai vektor tetap dan selisih
    vektor dengan citranya.
\end{latihan}

\begin{latihan}\label{op:lat-gram-sama}
    Misalkan $v_1,\ldots,v_m$ dan $w_1,\ldots,w_m$ dua daftar
    dalam ruang hasil kali dalam $V$ dengan
    $\langle v_i,v_j\rangle=\langle w_i,w_j\rangle$ untuk semua $i,j$.
    Buktikan bahwa ada operator uniter atau ortogonal $U$ dengan
    $Uv_j=w_j$ untuk setiap $j$. Daftar tidak harus bebas linear.
\end{latihan}

\begin{latihan}\label{op:lat-norma-sama}
    Untuk transformasi linear $A,B:V\to W$, andaikan
    $\|Ax\|=\|Bx\|$ untuk semua $x$.
    Buktikan bahwa ada operator uniter atau ortogonal $U$ pada $W$
    sehingga $A=UB$.
\end{latihan}

\begin{latihan}\label{op:lat-kuadratik-real}
    Pada ruang hasil kali dalam kompleks, buktikan bahwa
    $\langle x,Tx\rangle$ real untuk semua $x$ jika dan hanya jika
    $T$ adjoin diri. Jelaskan mengapa pernyataan ini tidak menjadi
    kriteria adjoin diri pada ruang real.
\end{latihan}

\begin{latihan}\label{op:lat-isotropik}
    Misalkan $T$ adjoin diri mempunyai $p$ nilai eigen positif,
    $q$ negatif, dan $z$ nol, dihitung dengan multiplisitas.
    Tentukan dimensi terbesar subruang $W$ yang memenuhi
    $\langle x,Tx\rangle=0$ untuk setiap $x\in W$.
\end{latihan}

\begin{latihan}\label{op:lat-pilih-metrik}
    Buktikan bahwa suatu operator pada ruang linear berdimensi
    hingga dapat dibuat adjoin diri dengan memilih suatu hasil
    kali dalam tepat ketika operator itu dapat didiagonalkan
    dengan semua nilai eigen real.
    Simpulkan bahwa setiap operator idempoten dapat dibuat
    menjadi proyeksi ortogonal untuk suatu hasil kali dalam.
\end{latihan}

\ifpembahasan
\section{Pembahasan Latihan}

\subsection*{Latihan \ref{op:lat-riesz}: evaluasi sebagai hasil kali dalam}
Tuliskan representer $r(t)=a+bt+ct^2$.
Syarat $\langle r,1\rangle=1$, $\langle r,t\rangle=0$,
dan $\langle r,t^2\rangle=0$ memberi
\[
    2a+\frac23c=1,\qquad \frac23b=0,\qquad
    \frac23a+\frac25c=0.
\]
Maka $r(t)=9/8-15t^2/8$.
Norma kuadrat representer adalah $\langle r,r\rangle=\ell(r)=9/8$.
Jadi konstanta terbaik $C=9/8$, dengan kesamaan untuk
polinom yang merupakan kelipatan $r$, termasuk polinom nol.

\subsection*{Latihan \ref{op:lat-adjoin-berbobot}: transpose saja belum cukup}
Rumus adjoin dalam basis ini memberi
\[
    [T^*]=G^{-1}A^TG
       =\frac13\begin{pmatrix}2&-1\\-1&2\end{pmatrix}
         \begin{pmatrix}2&1\\4&5\end{pmatrix}
       =\begin{pmatrix}0&-1\\2&3\end{pmatrix}.
\]
Matriks tersebut berbeda dari $A$, sehingga $T$ bukan adjoin diri.
Ia juga berbeda dari $A^T$, sebab basis standar tidak ortonormal
untuk metrik $G$.

\subsection*{Latihan \ref{op:lat-komutator}: selisih nilai eigen muncul berpasangan}
Rumus adjoin pada ruang matriks memberi $D_A^*=D_{A^*}=D_A$.
Pilih basis eigen ortonormal $u_1,\ldots,u_n$ bagi $A$,
dengan nilai eigen real $\lambda_j$.
Matriks $u_iu_j^*$ adalah vektor eigen $D_A$ untuk nilai
$\lambda_i-\lambda_j$.
Jika $D_A$ semidefinit positif, semua selisih itu nonnegatif.
Menukar $i,j$ memaksa semuanya nol, sehingga semua nilai
eigen $A$ sama. Teorema spektral memberi $A=cI$.
Sebaliknya, $A=cI$ menghasilkan $D_A=0$, yang semidefinit positif.

\subsection*{Latihan \ref{op:lat-titik-tetap}: kernel adjoin sama dengan kernel asal}
Karena $U^*=U^{-1}$, persamaan $U^*x=x$ ekuivalen dengan $Ux=x$.
Jadi $\ker((U-I)^*)=\ker(U-I)$.
Identitas kernel--citra memberi
\[
    \bigl(\operatorname{im}(U-I)\bigr)^\perp=\ker(U-I).
\]
Dekomposisi ortogonal mengikuti dengan mengambil komplemen.
Kernel berisi vektor yang tetap oleh $U$, sedangkan citra
berisi semua $Ux-x$. Mengganti tandanya memberi himpunan
yang sama, sehingga dapat pula ditulis sebagai semua $x-Ux$.

\subsection*{Latihan \ref{op:lat-gram-sama}: perluas isometri pada ruang pembangun}
Definisikan dahulu $U_0(\sum c_jv_j)=\sum c_jw_j$.
Kesamaan matriks Gram memberi
\[
    \left\|\sum c_jv_j\right\|^2
      =\left\|\sum c_jw_j\right\|^2.
\]
Jika dua representasi sumber sama, selisihnya bernorma nol;
persamaan tersebut memaksa selisih citranya nol juga.
Jadi $U_0$ terdefinisi dengan baik, linear, dan merupakan
isometri dari ruang pembangun $v_j$ ke ruang pembangun $w_j$.
Pilih basis ortonormal pada ruang pertama; citranya menjadi
basis ortonormal pada ruang kedua. Perluas kedua basis ke
basis ortonormal $V$, lalu kirim vektor tambahan secara berpasangan.
Operator yang dihasilkan uniter atau ortogonal dan memenuhi semua syarat.

\subsection*{Latihan \ref{op:lat-norma-sama}: cocokkan kedua citra}
Definisikan $U_0(Bx)=Ax$ pada $\operatorname{im}B$.
Jika $Bx=By$, hipotesis pada $x-y$ memberi
$\|A(x-y)\|=\|B(x-y)\|=0$, sehingga definisinya sah.
Pemetaan ini linear, mempertahankan norma, dan surjektif ke
$\operatorname{im}A$. Dengan polarisasi, ia mempertahankan
hasil kali dalam. Perluas isometri ini ke operator uniter
atau ortogonal pada $W$ dengan melengkapi basis ortonormal
kedua citra seperti pada latihan sebelumnya. Maka $U(Bx)=Ax$
untuk setiap $x$, yakni $A=UB$.

\subsection*{Latihan \ref{op:lat-kuadratik-real}: pisahkan dua bagian adjoin diri}
Arah perlu mengikuti simetri konjugat dan $T^*=T$.
Untuk arah cukup, tulis
\[
    T=H+iK,\qquad H=\frac{T+T^*}{2},\qquad
    K=\frac{T-T^*}{2i}.
\]
Keduanya adjoin diri. Karena
$\langle x,Tx\rangle=\langle x,Hx\rangle+i\langle x,Kx\rangle$
real, diperoleh $\langle x,Kx\rangle=0$ untuk semua $x$.
Uji pada basis eigen ortonormal $K$ menunjukkan bahwa setiap
nilai eigennya nol, sehingga $K=0$ dan $T=H$ adjoin diri.
Pada ruang real, setiap hasil kali dalam bernilai real secara
otomatis. Rotasi $\begin{pmatrix}0&-1\\1&0\end{pmatrix}$,
misalnya, tidak adjoin diri walaupun semua nilai kuadratnya real.

\subsection*{Latihan \ref{op:lat-isotropik}: batas dimensi dan konstruksi}
Tuliskan dekomposisi spektral ortogonal
$V=V_+\oplus V_-\oplus V_0$, dengan dimensi $p,q,z$.
Jika bentuk kuadratik nol pada $W$, maka
$W\cap V_+=W\cap V_-=\{0\}$.
Rumus dimensi memberi $\dim W\le q+z$ dan $\dim W\le p+z$.
Jadi batas atasnya $z+\min\{p,q\}$.
Untuk mencapai batas, pilih vektor eigen ortonormal $e_i$ untuk
nilai positif $\lambda_i$ dan $f_i$ untuk nilai negatif $-\mu_i$,
dengan $\lambda_i,\mu_i>0$.
Gabungkan basis $V_0$ dengan
\[
    w_i=\frac{e_i}{\sqrt{\lambda_i}}+
        \frac{f_i}{\sqrt{\mu_i}},\qquad 1\le i\le\min\{p,q\}.
\]
Daftar itu bebas linear. Untuk setiap kombinasinya, kontribusi
positif dan negatif pada bentuk kuadratik saling meniadakan,
sedangkan komponen $V_0$ tidak menyumbang apa pun.
Jadi dimensi terbesar tepat $z+\min\{p,q\}$.

\subsection*{Latihan \ref{op:lat-pilih-metrik}: jadikan basis eigen ortonormal}
Jika operator adjoin diri untuk suatu hasil kali dalam,
teorema spektral memberi basis eigen dan nilai eigen real.
Sebaliknya, pilih basis eigen dengan nilai eigen real dan
definisikan hasil kali dalam dengan menjadikan basis itu ortonormal.
Matriks operator dalam basis tersebut diagonal real, sehingga
operator adjoin diri.
Jika matriks kolom basis eigen dalam koordinat awal adalah $S$,
pilihan metriknya secara eksplisit
$G=(S^{-1})^*S^{-1}$.
Operator idempoten mempunyai dekomposisi
$V=\operatorname{im}P\oplus\ker P$, dan matriksnya dalam basis
yang sesuai adalah $\diag(I,0)$.
Jadikan basis tersebut ortonormal. Citra dan kernel sekarang
ortogonal, sehingga $P$ menjadi proyeksi ortogonal.
\fi
